\documentclass[11pt]{article}
\setlength{\textwidth}{6.6in}
\setlength{\textheight}{9.35in}
\setlength{\oddsidemargin}{-0.05in}
\setlength{\evensidemargin}{-0.05in}
\setlength{\topmargin}{-0.65in}
\usepackage{amsmath,amssymb}
\usepackage{booktabs}
\usepackage{xcolor}

\definecolor{NavyBlue}{RGB}{0,51,102}
\definecolor{DarkGreen}{RGB}{0,100,0}

\newcommand{\commentbox}[2]{%
\vspace{0.42em}\noindent
\textbf{\textcolor{NavyBlue}{#1:}} \textit{#2}\vspace{0.18em}%
}
\newcommand{\responsebox}[1]{%
\noindent
\textbf{\textcolor{DarkGreen}{Response:}} #1\vspace{0.42em}%
}

\begin{document}
\thispagestyle{empty}

\begin{center}
{\Large \textbf{Point-by-Point Response to Reviewers and Editor}}\\[0.6em]
{\large Document Type: Point-by-Point Summary of Revisions Following Prior Peer Review}\\[0.4em]
\textbf{``Calibration-Aware Anisotropic 3D Hysteresis and Leakage Auditing for Volumetric Lower-Grade Glioma Segmentation and Radiogenomic Profiling in Brain MRI''}\\[0.4em]
Shayan Ganji Nataj, Keivan Borna (Corresponding Author)
\end{center}

\vspace{1.0em}
\noindent
This document summarizes our point-by-point responses to the methodological and clinical critiques raised during prior peer review of an earlier version of this study (Ref. IMAVIS-D-26-02558). We sincerely thank the previous Associate Editor and Reviewer for their constructive feedback. In this thoroughly revised manuscript submitted to Biomedical Signal Processing and Control, every identified issue has been fully resolved through strict nested cross-validation, controlled empirical leakage audits, probability calibration, and comprehensive 3D volumetric and radiogenomic analyses.

\section*{Response to the Associate Editor}

\commentbox{Associate Editor Comment 1}{The authors must ensure complete methodological rigor, eliminating any possibility of threshold tuning leakage or ambiguous ensembling, and clearly demonstrate the clinical and 3D utility of the proposed segmentation system.}
\responsebox{We deeply appreciate the Editor's guidance. In this revised submission, we have addressed all methodological concerns:
\begin{enumerate}
    \item \textbf{Strict Nested Cross-Validation:} All hyperparameter tuning---including temperature scaling calibration ($T_k^*$), horizontal in-plane binarization ($\tau_{xy, k}^*$), axial rescue threshold ($\tau_{z,\mathrm{low}, k}^*$), and 3D component filtering ($V_{min, k}^*, \tau_{seed, k}^*$)---is performed exclusively on the internal validation set (\texttt{inner\_val}) of each fold. The test set (\texttt{outer\_test}) remains untouched until final evaluation. We formally quantify the artificial optimism of post-hoc test tuning as $\Delta\text{Optimism} = +0.063\%$ in Section~5.3.
    \item \textbf{Unambiguous Evaluation Mapping:} We explicitly clarify that out-of-fold (OOF) predictions are generated strictly by the dedicated model trained for that fold (a 1-to-1 patient mapping), completely refuting any 5-model ensemble averaging ambiguity (Section~3.3).
    \item \textbf{Clinical and 3D Volumetric Agreement:} We reconstruct full 3D volumes for all 110 TCGA-LGG patients, evaluating 3D Volumetric Dice ($77.73\%$), 95\% Hausdorff Distance ($10.97$ voxels), Average Surface Distance ($3.10$ voxels), Intraclass Correlation ($\text{ICC}(2, 1) = 0.8894$), and genomic subtype concordance ($p = 0.0268$ across \textit{COCCluster} cohorts).
\end{enumerate}
}

\section*{Response to Reviewer 1}

\commentbox{Comment 1 (Threshold Optimization and Optimism Bias)}{The selection of the decision threshold appears to have been optimized on the testing data, which introduces data leakage and optimistic performance reporting. The authors should use a nested validation framework or report fixed-threshold baselines.}
\responsebox{We fully agree with the Reviewer. In the revised manuscript, we implemented a strict nested 5-fold cross-validation scheme:
\begin{itemize}
    \item In each fold $k$, the optimal threshold $\tau_{xy, k}^*$ is selected solely by maximizing positive-slice Dice on the unseen \texttt{inner\_val} partition.
    \item The test partition is evaluated exclusively using the frozen $\tau_{xy, k}^*$.
    \item To directly address the Reviewer's concern, Section~5.3 and Table~4 present a formal comparison between:
    \begin{enumerate}
        \item Fixed Baseline ($\tau=0.50$): $\text{Dice}_{\text{pos}} = 73.96\%$ [95\% CI: 70.10\%--77.15\%]
        \item Nested Inner-Fold Search ($\tau_k^*$): $\text{Dice}_{\text{pos}} = 74.60\%$ [95\% CI: 70.45\%--78.08\%]
        \item Leaked Pooled Search ($\tau_{\text{leak}}^* = 0.88$ on test data): $\text{Dice}_{\text{pos}} = 74.66\%$
    \end{enumerate}
    \item We formally define and quantify the \textbf{Optimism Bias}:
    \begin{equation*}
    \Delta\text{Optimism} = \text{Dice}(\tau_{\text{leak}}^*) - \text{Dice}(\tau_{xy, k}^*) = 74.66\% - 74.60\% = +0.063\%
    \end{equation*}
    This demonstrates that our reported gains are nested inner-fold calibrated without test snooping.
\end{itemize}
}

\commentbox{Comment 2 (Ensemble Ambiguity in OOF Predictions)}{It is unclear whether out-of-fold (OOF) results were produced by averaging predictions across all 5 fold models (an ensemble) or by evaluating each test fold independently on its corresponding trained model.}
\responsebox{We thank the Reviewer for highlighting this ambiguity. We have rewritten Section~3.3 to state the mathematical mapping unambiguously:
\begin{equation}
\hat{Y}_i = \mathcal{M}_k(X_i; \theta_k^*), \quad \forall i \in \mathcal{D}_{\text{test}}^{(k)}, \quad k \in \{1, 2, 3, 4, 5\}
\end{equation}
Every patient slice belongs to exactly one test fold $\mathcal{D}_{\text{test}}^{(k)}$ and is segmented solely by the single network $\mathcal{M}_k$ trained on folds $j \neq k$. No model averaging, soft blending, or test-time ensembling was performed. The resulting 3,929 out-of-fold predictions represent true prospective generalization to unseen patients.}

\commentbox{Comment 3 (Slice Leakage vs. Patient-Grouped Splitting)}{Many previous 2D MRI segmentation papers split datasets at the slice level rather than the patient level, leading to memorization of contiguous slices. The authors should evaluate whether slice-level splitting significantly inflates results.}
\responsebox{The Reviewer raises a foundational concern in medical image analysis. To answer this definitively, we conducted an aligned, matched 5-fold cross-validation experiment comparing:
\begin{itemize}
    \item \textbf{Patient-Grouped Split (Zero Leakage):} All slices belonging to any given patient are assigned strictly to a single fold.
    \item \textbf{Random-Slice Split (Data-Leaked):} Slices from the same patient are randomly dispersed across training and testing partitions.
\end{itemize}
As reported in Section~5.5 and Table~6, the random slice split yields an artificial positive-slice Dice inflation of \textbf{+8.79\%} (82.74\% vs. 73.96\%) and an overall Dice inflation of \textbf{+6.20\%} (88.22\% vs. 82.02\%). Furthermore, Figure~4(c)--(d) illustrates the severe curve inversion in the leaky split, where validation Dice abnormally exceeds training Dice due to slice memorization. This controlled experiment proves that high Dice numbers in prior 2D literature often stem from slice leakage rather than architectural innovation.}

\commentbox{Comment 4 (Validation Measurement Protocol and Data Augmentation)}{Were validation loss and Dice curves computed with stochastic data augmentation enabled or in pure inference mode? Augmentation during validation distorts performance curves.}
\responsebox{In our revised pipeline (Section~3.2), validation and test evaluation are conducted in \textbf{pure inference mode} without data augmentation. Stochastic transformations (affine rotations, scaling, and elastic deformation) are active strictly during forward-backward training steps. Slices in \texttt{inner\_val} and \texttt{outer\_test} are fed in deterministic axial sequence using preprocessed, unaugmented MRI tensors. The convergence history across all 5 folds is illustrated in Figure~4(a)--(b).}

\commentbox{Comment 5 (Artifact Screening Audit and Slices Removed)}{The manuscript mentions excluding 200 artifact slices during training. The authors must audit whether this exclusion introduces selection bias or discards small neoplastic lesions.}
\responsebox{We have conducted a slice-by-slice morphological and histopathological audit across all 3,929 slices (Section~2.2, Table~2, and Figure~1):
\begin{enumerate}
    \item \textbf{Empty Slices:} Out of the 200 excluded slices, exactly \textbf{199 slices (99.50\%)} are completely tumor-free ($y = 0$). Only 1 slice contained a marginal lesion (0.50\%, Slice 23 of TCGA\_DU\_5855) where severe gradient coil clipping truncated the brain parenchyma.
    \item \textbf{Zero Change on Positive Slices:} Evaluating the model on the full cohort (3,929 slices) versus the screened cohort (3,729 slices) results in \textbf{identical positive-slice Dice of 73.96\%} (a change of exactly $0.00\%$, Table~2).
    \item \textbf{Wavelet Denoising Gate:} Denoising is selectively applied only when slice quality index $Q(I) < 4.0$ (3,738 slices, 95.14\%), while high-quality slices ($Q \ge 4.0$, 191 slices, 4.86\%) bypass filtering to preserve fine parenchymal microstructures (illustrated in Figure~1).
\end{enumerate}
}

\commentbox{Comment 6 (Calibration and Uncertainty Estimation)}{Deep segmentation models trained with overlap-based loss functions (such as Focal Tversky) tend to be severely uncalibrated and overconfident. How does the proposed method address probabilistic reliability?}
\responsebox{We thoroughly investigated probability calibration in Section~4.1 and Section~5.2. In our Focal Tversky loss formulation:
\begin{equation}
\text{TI}(y, \hat{p}) = \frac{\sum_i y_i \hat{p}_i + \epsilon}{\sum_i y_i \hat{p}_i + \alpha \sum_i y_i(1 - \hat{p}_i) + \beta \sum_i (1 - y_i)\hat{p}_i + \epsilon}
\end{equation}
the coefficient $\alpha = 0.80$ is the weight penalizing \textbf{False Negatives (FN)}, while $\beta = 0.20$ penalizes \textbf{False Positives (FP)}. By heavily penalizing missed lesions ($\alpha = 0.80 \gg \beta = 0.20$), the network prioritizes high sensitivity/Recall ($82.98\%$) to capture delicate infiltrative boundaries. However, this asymmetric penalty causes the raw network to emit an overconfident false-positive halo around lesions (Precision = $71.95\%$) and severe overconfidence ($T^* = 1.844$). 

To restore probabilistic reliability without altering ROC discrimination, we applied nested scalar temperature scaling:
\begin{equation}
\hat{P}_{i, z} = \sigma\left(\frac{z_{i}}{\hat{T}_k^*}\right), \quad \hat{T}_k^* = \arg\min_{T > 0} \mathcal{L}_{\text{NLL}}\left(\frac{z_{\text{val}, k}}{T}, y_{\text{val}, k}\right)
\end{equation}
Across the 5 folds, the optimal temperature averaged $\hat{T}^* = 1.844$, confirming that raw logits were overconfident ($T > 1$). On untouched test data:
\begin{itemize}
    \item Pooled cohort Expected Calibration Error (ECE) dropped by \textbf{46.9\%} (0.369\% to 0.196\%, Table~4 and Table~6).
    \item Mean per-fold ECE dropped by \textbf{31.7\%} (0.372\% to 0.254\%).
    \item Brier Score improved from 0.00375 to 0.00338 (-9.9\%), and NLL decreased from 0.0230 to 0.0154 (-32.9\%).
    \item Reliability diagrams (Figure~3) confirm that empirical probability bins align closely with the diagonal identity line after temperature scaling.
\end{itemize}
Subsequently, our nested thresholding ($\tau_{xy}^* = 0.74$) and Anisotropic 3D Z-Hysteresis shave off this FP halo, elevating Precision to $76.86\%$ (Table~5).}

\commentbox{Comment 7 (Handling of Non-Tumor Slices in Evaluation)}{In brain MRI, many slices contain no tumor. How are empty slices defined and penalized in Dice calculations, and what is the false-positive rate on healthy slices?}
\responsebox{In Section~3.4, we explicitly state our mathematical convention for slice-level metrics:
\begin{equation}
\text{Dice}(Y, \hat{Y}) = \begin{cases}
1.0, & \text{if } |Y| = 0 \text{ and } |\hat{Y}| = 0 \text{ (True Negative)} \\
0.0, & \text{if } |Y| = 0 \text{ and } |\hat{Y}| > 0 \text{ (False Alarm)} \\
\frac{2|Y \cap \hat{Y}|}{|Y| + |\hat{Y}|}, & \text{if } |Y| > 0
\end{cases}
\end{equation}
In the TCGA-LGG cohort, 2,556 out of 3,929 slices (65.05\%) are completely healthy. To prevent metric distortion, we report:
\begin{itemize}
    \item $\text{Dice}_{\text{pos}}$: Evaluated strictly on the 1,373 tumor-bearing slices ($74.92\%$).
    \item $\text{Dice}_{\text{all}}$: Evaluated across all 3,929 slices including true negatives ($86.12\%$).
    \item $\text{FPR}_{\text{empty}}$: False alarm rate on empty slices.
\end{itemize}
Crucially, our Proposed Anisotropic 3D Z-Hysteresis reduces $\text{FPR}_{\text{empty}}$ from \textbf{13.65\% to 7.86\%} (a \textbf{42.4\% relative reduction} in false alarm slices, Table~5), while boosting $\text{Dice}_{\text{all}}$ from 82.02\% to \textbf{86.12\%}.}

\commentbox{Comment 8 (Fair Comparison with State-of-the-Art)}{The authors should provide a comprehensive and fair comparison with previous literature on the TCGA-LGG benchmark, detailing splitting regimes, post-processing methods, and model complexities.}
\responsebox{In Section~7.1 and Table~8, we present a rigorous benchmark comparison against genuine published studies evaluated strictly on the Buda et al. TCGA-LGG cohort (Comput. Biol. Med. 2019). We explicitly categorize studies by:
\begin{enumerate}
    \item \textbf{Split Protocol:} Patient-Grouped (Zero Leakage) vs. Random-Slice (Data-Leaked).
    \item \textbf{Metric Formulation:} Thresholded hard $\text{Dice}_{\text{pos}}$ vs. unthresholded continuous soft Dice.
    \item \textbf{Calibration \& 3D Post-Processing:} Temperature scaling and Anisotropic 3D Z-Hysteresis vs. uncalibrated raw argmax.
    \item \textbf{Parameter Footprint:} Compact models (Yuan \& Kang 2025: 1.61M; Ours: 8.20M) vs. standard models (Buda 2019: 7.80M).
\end{enumerate}
Under strictly patient-isolated cross-validation with nested calibration, our framework establishes the leading performance on TCGA-LGG ($\text{Dice}_{\text{pos}} = 74.92\%$, $\text{Dice}_{\text{all}} = 86.12\%$, 3D Volumetric $\text{Dice} = 77.73\%$, $\text{ICC} = 0.8894$) while operating with an efficient 8.20M parameter footprint.}

\vspace{1.5em}
\noindent
We believe that these extensive revisions, rigorous audits, and new clinical experiments fully address all comments and significantly elevate the scientific contribution of our manuscript.

\end{document}
